\documentclass[11pt,a4paper]{article}

\usepackage[utf8]{inputenc}
\usepackage[T1]{fontenc}
\usepackage[spanish,es-tabla]{babel}
\usepackage[margin=2.5cm]{geometry}
\usepackage{array}
\usepackage{longtable}
\usepackage{hyperref}

\newcolumntype{L}[1]{>{\raggedright\arraybackslash}p{#1}}

\title{FORM-TL-01 -- Taller de tipos de requisitos\\
\large Ingeniería de Requerimientos (ISR-401) -- 4.\textsuperscript{o} nivel -- Paralelo A}
\author{Alan Andrés Aldas Zambrano}
\date{Fecha límite de entrega: 7 de noviembre de 2026}

\begin{document}

\maketitle

% ---------------------------------------------------------------
\section{Clasificación de los enunciados}

% Sistema: Plataforma de gestión de citas médicas. Enunciados copiados de la Tabla 1 del caso entregado por el docente.
% Tipo: funcional, calidad (indicar la característica de ISO/IEC 25010:2023) o restricción.

\begin{longtable}{|c|L{3.6cm}|L{2.4cm}|L{1.6cm}|L{1.6cm}|L{3.4cm}|}
\caption{Clasificación de los 20 enunciados.}\label{tab:clasificacion}\\
\hline
\textbf{Código} & \textbf{Enunciado} & \textbf{Tipo} & \textbf{Producto o proceso} & \textbf{Sistema o software} & \textbf{Justificación} \\
\hline
\endfirsthead
\hline
\textbf{Código} & \textbf{Enunciado} & \textbf{Tipo} & \textbf{Producto o proceso} & \textbf{Sistema o software} & \textbf{Justificación} \\
\hline
\endhead
E01 & El paciente podrá reservar una cita eligiendo especialidad, médico, fecha y hora entre los horarios disponibles. & & & & \\ \hline
E02 & La plataforma enviará al paciente un recordatorio por mensaje de texto 24 horas antes de su cita. & & & & \\ \hline
E03 & La aplicación debe ser rápida. & & & & \\ \hline
E04 & El recepcionista podrá registrar la llegada del paciente que se acerca a la ventanilla. & & & & \\ \hline
E05 & La plataforma se desarrollará en Java 21 con base de datos PostgreSQL, tecnologías que domina el área de TI. & & & & \\ \hline
E06 & El médico podrá consultar la lista de pacientes del día y marcar cada consulta como atendida. & & & & \\ \hline
E07 & Los datos de salud de los pacientes deben estar seguros. & & & & \\ \hline
E08 & La aplicación móvil debe funcionar en teléfonos con Android 10 o superior. & & & & \\ \hline
E09 & Una cita cancelada con menos de cuatro horas de anticipación se registrará como inasistencia. & & & & \\ \hline
E10 & La plataforma debe ser fácil de usar para los pacientes adultos mayores. & & & & \\ \hline
E11 & El proyecto debe entregarse en un plazo máximo de seis meses. & & & & \\ \hline
E12 & La sesión de un usuario se cerrará automáticamente tras 10 minutos sin actividad. & & & & \\ \hline
E13 & El quiosco de la sala de espera leerá el número de cédula ingresado en su pantalla táctil e imprimirá un ticket con el número de turno. & & & & \\ \hline
E14 & La plataforma debe estar siempre disponible. & & & & \\ \hline
E15 & El director médico podrá generar un reporte mensual de inasistencias por especialidad. & & & & \\ \hline
E16 & Cada incremento se validará con el director médico al final de cada iteración de dos semanas. & & & & \\ \hline
E17 & El tratamiento de los datos personales de los pacientes debe cumplir la Ley Orgánica de Protección de Datos Personales del Ecuador. & & & & \\ \hline
E18 & La plataforma debe integrarse sin problemas con el sistema de facturación actual. & & & & \\ \hline
E19 & Si el servidor principal falla, la plataforma debe restablecer el servicio en menos de 30 minutos sin perder ninguna cita confirmada. & & & & \\ \hline
E20 & Todo cambio en el código fuente será revisado por otro integrante del equipo antes de integrarse a la rama principal. & & & & \\ \hline
\end{longtable}

% ---------------------------------------------------------------
\section{Requisitos de calidad cuantificados}

\begin{longtable}{|c|L{2.8cm}|L{5.2cm}|L{2.4cm}|L{2.4cm}|}
\caption{Reescritura de cinco requisitos de calidad vagos.}\label{tab:cuantificados}\\
\hline
\textbf{N.\textsuperscript{o}} & \textbf{Enunciado vago} & \textbf{Requisito cuantificado} & \textbf{Característica ISO/IEC 25010} & \textbf{Métrica y umbral} \\
\hline
\endfirsthead
\hline
\textbf{N.\textsuperscript{o}} & \textbf{Enunciado vago} & \textbf{Requisito cuantificado} & \textbf{Característica ISO/IEC 25010} & \textbf{Métrica y umbral} \\
\hline
\endhead
1 & & & & \\ \hline
2 & & & & \\ \hline
3 & & & & \\ \hline
4 & & & & \\ \hline
5 & & & & \\ \hline
\end{longtable}

% ---------------------------------------------------------------
\section{Discrepancias con el compañero}

Compañero con quien se comparó: % escribir el nombre

\begin{longtable}{|c|L{3.2cm}|L{3.2cm}|L{5.4cm}|}
\caption{Discrepancias encontradas y su resolución.}\label{tab:discrepancias}\\
\hline
\textbf{Enunciado} & \textbf{Mi clasificación} & \textbf{Clasificación del compañero} & \textbf{Resolución y argumento} \\
\hline
\endfirsthead
\hline
\textbf{Enunciado} & \textbf{Mi clasificación} & \textbf{Clasificación del compañero} & \textbf{Resolución y argumento} \\
\hline
\endhead
 & & & \\ \hline
\end{longtable}

% ---------------------------------------------------------------
\section{Uso de IA como apoyo}

% Completar con la verdad: herramienta, versión, propósito y cómo se verificó lo obtenido.

\begin{itemize}
  \item \textbf{Herramienta y versión:}
  \item \textbf{Propósito:}
  \item \textbf{Verificación:}
\end{itemize}

% ---------------------------------------------------------------
% \nocite{*} muestra todas las referencias del .bib aunque todavía no se citen.
% Cuando cites en el texto con \cite{pohl2025}, \cite{iso29148} o \cite{iso25010}, puedes quitarlo.
\nocite{*}
\bibliographystyle{IEEEtran}
\bibliography{referencias}

\end{document}
